\chapter{Earache}
\index{Earache}

\section*{Definition and Overview}
Earache is pain in or around the ear. It is one of the most common complaints in both children and adults and has many causes, ranging from minor problems such as pressure changes and earwax to infections of the ear canal, middle ear, and throat. Pain that is felt in the ear but originates elsewhere, such as from the teeth or jaw, is called referred pain. This chapter explains the common causes of earache, how to tell when it is serious, and how it is treated, including when antibiotics are needed and when they are not.

\section*{Epidemiology}
Earache is extremely common, particularly in children, and is one of the leading reasons for GP and emergency department visits in the early years of life. Middle ear infections are most common in children under three, while outer ear infections are common in swimmers and adults. Many episodes of earache, especially in older children and adults, are caused by viral infections or eustachian tube pressure and settle without treatment.

\section*{Aetiology and Pathophysiology}
The causes of earache can be divided into those affecting the ear itself and those causing referred pain.
\begin{itemize}
    \item \textbf{Middle ear infection (otitis media):} infection behind the eardrum, usually after a cold, causing pressure and pain
    \item \textbf{Outer ear infection (otitis externa):} infection of the ear canal, common after swimming or injury
    \item \textbf{Eustachian tube dysfunction:} pressure changes with flying, altitude, or colds, when the tube connecting the middle ear to the throat fails to equalise pressure
    \item \textbf{Earwax impaction:} wax pressing on the eardrum
    \item \textbf{Foreign bodies:} objects inserted into the ear, particularly in children
    \item \textbf{Injury:} trauma to the ear, including perforation of the eardrum
    \item \textbf{Referred pain:} from dental problems, tonsillitis, throat infections, and jaw joint disorders
    \item \textbf{Other causes:} eczema or psoriasis of the canal, boils, and rarely tumours
\end{itemize}

\section*{Clinical Features}
The character of the pain and the associated symptoms point to the cause.
\begin{itemize}
    \item Pain in, behind, or around the ear
    \item Pain that is worse when lying down, chewing, or pulling on the ear
    \item Fever, irritability, and pulling at the ear in children with middle ear infection
    \item Itching, discharge, and pain on touching the ear in outer ear infection
    \item A feeling of fullness, popping, or crackling with eustachian tube problems
    \item Muffled hearing
    \item Tinnitus or dizziness
    \item Sore throat, toothache, or jaw pain suggesting referred pain
\end{itemize}

\section*{Differential Diagnosis}
Earache can be confused with, or caused by, other conditions.
\begin{itemize}
    \item Dental problems, including cavities, abscesses, and erupting wisdom teeth
    \item Temporomandibular joint disorders and teeth grinding
    \item Tonsillitis, pharyngitis, and throat infections
    \item Sinusitis
    \item Neck and cervical spine problems
    \item Trigeminal neuralgia and other nerve pain
    \item Headache disorders
\end{itemize}

\section*{When to Seek Medical Care}
See a doctor for earache that is severe, persistent, or accompanied by fever, discharge, or hearing loss, and for any earache in an infant or young child.

\textbf{Contact your local emergency services for:}
\begin{itemize}
    \item Severe pain with swelling and redness behind the ear
    \item Facial weakness or drooping
    \item Stiff neck, confusion, or high fever, which may indicate spread of infection
    \item Fluid or blood draining from the ear after a head injury
    \item Sudden severe headache or seizures
\end{itemize}

\section*{Diagnosis and Investigations}
Most earache is diagnosed clinically.
\begin{itemize}
    \item \textbf{History:} the nature, onset, and context of the pain, and any associated symptoms
    \item \textbf{Otoscopy:} examination of the ear canal and eardrum
    \item \textbf{Pneumatic otoscopy:} a puff of air to check eardrum movement
    \item \textbf{Tympanometry:} assessment of middle ear pressure
    \item \textbf{Hearing tests:} when hearing loss is present
    \item \textbf{Dental and throat examination:} to identify referred causes
\end{itemize}

\section*{Management}
Treatment depends on the cause.
\begin{itemize}
    \item \textbf{Pain relief:} paracetamol or ibuprofen is the most important treatment for most earache
    \item \textbf{Watchful waiting:} many middle ear infections settle within days without antibiotics
    \item \textbf{Antibiotics:} for severe or persistent middle ear infections, children under two, or those at risk of complications
    \item \textbf{Ear drops:} for outer ear infections
    \item \textbf{Wax removal:} for impacted wax
    \item \textbf{Pressure equalisation:} swallowing, yawning, or special techniques for flying-related earache
    \item \textbf{Treatment of the cause:} dental treatment, throat care, and management of jaw problems for referred pain
\end{itemize}

\section*{Complications}
\begin{itemize}
    \item Perforation of the eardrum
    \item Mastoiditis, an infection of the bone behind the ear
    \item Chronic middle ear infection and glue ear, affecting hearing
    \item Spread of infection in rare cases
    \item Persistent pain and anxiety if the cause is not identified
\end{itemize}

\section*{Prognosis and Outcomes}
Most earache resolves within a few days, and the outlook is excellent with appropriate treatment of the cause. Middle ear infections in children settle without long-term effects in the vast majority of cases, and hearing returns to normal once fluid clears. Prompt attention to severe or persistent pain, and to pain in young children, prevents complications.

\section*{Prevention and Self-Care}
\begin{itemize}
    \item Relieve pressure during flights by swallowing, yawning, or chewing
    \item Keep children's immunisations up to date
    \item Avoid exposing children to cigarette smoke
    \item Dry ears after swimming and avoid inserting objects into the ear
    \item Use paracetamol or ibuprofen for pain, following the dose for age
    \item See a doctor if the pain persists, worsens, or is severe
    \item Keep the ear dry if there is discharge
\end{itemize}

\section*{Sources and Further Reading}
The following reputable sources provide further information for healthcare students and patients seeking reliable, current advice about earache.
\begin{itemize}
    \item Healthdirect Australia. \textit{Earache: causes and treatment.} https://www.healthdirect.gov.au/earache
    \item NHS. \textit{Earache: causes and treatment.} https://www.nhs.uk/conditions/earache/
    \item Mayo Clinic. \textit{Ear pain: causes and treatment.} https://www.mayoclinic.org/diseases-conditions/ear-pain/
    \item Royal Children's Hospital Melbourne. \textit{Earache and ear infection fact sheet.} https://www.rch.org.au/
    \item Better Health Channel. \textit{Earache and ear problems.} https://www.betterhealth.vic.gov.au/
    \item MSD Manual. \textit{Earache: patient overview.} https://www.msdmanuals.com/
\end{itemize}
